%%%%%%%%%%%%%%%%%%%%%%%%%%%%%%%%%%%%%%%%%%%%%%%%%%%%%%%%%%%%%%%%%%%%%
%% sec_theory_EM_v2.tex
%%
%% Sections 2.2 -- 2.4 for paperH_HLmain.tex
%% RMG Periodic RT-TDDFT -- Paper 2
%% REVISED April 2026: clean separation of EM-field perturbations
%% (gauge-dependent, Track A) from scalar-potential perturbations
%% (gauge-independent, Track B).
%%
%% Subsection structure:
%%   2.2  Interaction with electromagnetic radiation
%%        2.2.1  Plane wave and minimal coupling
%%        2.2.2  Multipole expansion and long-wavelength approximation
%%        2.2.3  Crystal momentum conservation: vertical transitions
%%        2.2.4  Beyond the optical limit: scalar-potential perturbations
%%   2.3  Gauge transformation: length to velocity gauge
%%        2.3.1  Why the length gauge fails for periodic systems
%%        2.3.2  The Goppert-Mayer transformation
%%        2.3.3  The velocity-gauge TDKS Hamiltonian
%%        2.3.4  The perturbative momentum kick
%%   2.4  Vector potential and magnetic perturbations
%%        2.4.1  Electric perturbation: spatially uniform A(t)
%%        2.4.2  Magnetic perturbation: spatially varying A(r)
%%        2.4.3  Scope of the present implementation
%%
%% CONVENTIONS (all from notation.tex):
%%   Atomic units: hbar = m_e = e = 4*pi*eps_0 = 1
%%   Electron charge = -1  (e > 0 is elementary charge magnitude)
%%   Imaginary unit: \imath  (dotless i, consistent with paper 1)
%%   \myblue{} for all annotations and placeholders -- NEVER 
%%
%% SOURCES:
%%   NEAT-notes-RMG_main.tex lines 116-716 (primary)
%%   NEAT-project-notes_main.tex lines 256-716 (supplementary)
%%   Physics discussion notes, April 2026
%%%%%%%%%%%%%%%%%%%%%%%%%%%%%%%%%%%%%%%%%%%%%%%%%%%%%%%%%%%%%%%%%%%%%

%In this section we use Gaussian units to keep physical dimensions and the electron charge sign explicit throughout the derivations.

%\myblue{Move to IMPLEMENTATION: For the remainder of the paper we adopt atomic units 
%($\hbar=m_e=e= 4\pi\varepsilon_0 =1$ , $c= 1/\alpha \approx 137$)}
%(ℏ = mₑ = e = 4πε₀ = 1, c = 1/α ≈ 137). The key substitutions are: e/mc → 1/c = α, ℏ → 1, m → 1.}
%The fine-structure constant:
%α=(e^2)/(4 π ε0 ℏc)= e^2/ℏc(Gaussian)≈1/137.036
%\alpha = \frac{e^2}{4\pi\varepsilon_0\hbar c} = \frac{e^2}{\hbar c} \quad\text{(Gaussian)} \approx \frac{1}{137.036}

%In this section we discuss theory behind the current  RT-TDDFT implementation,

This section presents the theoretical foundations of real-time time-dependent density functional theory (RT-TDDFT)
for periodic systems interacting with electromagnetic radiation. We begin with the working equations for 
the electromagnetic field and the minimal coupling Hamiltonian (Section 2.1), then develop the gauge transformation 
from the length gauge to the velocity gauge (Section 2.2), the electronic current density operator including nonlocal 
pseudopotential contributions (Section 2.3), and the time propagation of the Kohn-Sham equations (Section 2.4). 
Rather than stating final expressions alone, we provide complete, step-by-step derivations with all intermediate 
results and sign conventions made explicit. 
Our motivation for this level of detail is practical: the existing literature on velocity-gauge RT-TDDFT is spread 
across several communities - solid-state physics, quantum chemistry, and optical response theory - each employing 
different conventions for the electron charge sign, the minimal coupling substitution, and the definition of the current operator. 
Reconstructing a consistent set of working equations from these sources requires reconciling notational conflicts 
and filling in steps that are frequently omitted. 
We aim to present a self-contained derivation that can serve as a reference for implementation. 
To this end, we work in Gaussian units throughout the theory section, keeping $\hbar$, $m$, and $e$ explicit 
so that physical dimensions and charge signs remain traceable at every step. 
Atomic units ($\hbar=m_e=e=4\pi\varepsilon_0=1)$ are adopted starting in Section 3 for the working equations as implemented in the code.


%--------------------------------------------------------------------
\subsection{Interaction of Matter with Electromagnetic Radiation}
\label{sec:em-interaction}
%--------------------------------------------------------------------

%The starting point for any RT-TDDFT calculation is the coupling of
%the electronic system to an external time-dependent perturbation.
%
%\subsection{External electromagnetic fields in periodic RT-TDDFT}
%\label{sec:external_fields}



The starting point for any RT-TDDFT calculation  is the time-dependent Kohn–Sham equation for single-particle orbitals 
$\Psi_j(r,t)$ in the presence of external electromagnetic fields which,  with  the minimal coupling framework, 
takes the form
\begin{equation}
	\imath \hbar \pdv{}{t} \Psi_j(\vb r,t) = \left[  \frac{1}{2m} \left( \vb{\hat p} - \frac{q_e}{c}\vb A_\mathrm{ext}(\vb r, t)\right)^2
	 + q_e \Phi_\mathrm{ext} (\vb r,t)   + \hat V_\mathrm{KS} [\rho](t)
	 \right] \Psi_j(\vb r,t)
\end{equation}
where $ \vb A_\mathrm{ext}$ and  $\Phi_\mathrm{ext}$ are vector and scalar potential representing external electromagnetic (EM) field, 
$\vb{\hat p} =-\imath\hbar \nabla $ is a canonical momentum operator,  
and $q_e=-1$ is and electron charge.
The Kohn-Sham potential, $ \hat V_\mathrm{KS} [\rho](t)$, 
 contains the ionic potential, Hartree potential, and exchange-correlation potential: 
\begin{equation}
  \hat V_{\mathrm{KS}} =   
  \hat V_{\mathrm{ion}} + \hat V_{\mathrm{H}}[\rho](\mathbf r,t) + \hat V_{\mathrm{xc}}[\rho](\mathbf r,t).
\end{equation}
In practical implementations based on norm-conserving or ultrasoft pseudopotentials,  $\hat V_{\mathrm{ion}}$ may contain both local
and nonlocal parts, and gauge transformations of the Hamiltonian may introduce additional terms that must be accounted for in both the
Hamiltonian and the current density operator.
The choice of the specific form of $\vb A_\mathrm{ext}$ and  $\Phi_\mathrm{ext}$ depends on the problem and is subject to a gauge freedom.\cite{gauge-theory}



 For a monochromatic external EM field, the vector potential may be represented as a plane wave,
\begin{equation}
	%\vb A(r,t)= \vb  A_0(\vb r,t) e^{\imath (\vb q \vb r -\omega t)} +  A_0^*( \vb r,t) e^{-\imath (\vb q \vb r -\omega t)}
	\vb A_\mathrm{ext}(r,t)= \vb  A_0 e^{\imath (\vb q \cdot \vb r -\omega t)} +  c.c. 
	\label{eq:plane_wave_vector_potential}
\end{equation}
where  $\omega$  is the angular frequency, a wave vector $\vb q$ describes propagation direction, and $\vb A_0=[A^0_x,A^0_y,A^0_z]$ is a complex polarization  amplitude.
The corresponding electric and magnetic fields are obtained from Maxwell equations:
\begin{align*}
	\mathbf E(\mathbf r,t) =& -\nabla \Phi_{\mathrm{ext}}(\mathbf r,t) - \frac{1}{c}\frac{\partial \mathbf A_{\mathrm{ext}}(\mathbf r,t)}{\partial t}, \\
	\mathbf B(\mathbf r,t) =& \nabla \times \mathbf A_{\mathrm{ext}}(\mathbf r,t).
\label{eq:em_fields_from_potentials}
\end{align*}
%Imposing the Coulomb gauge condition,
%leads to the transversality condition
%\begin{equation}
%  \vb q \cdot \vb A_0 = 0,
%\label{eq:transversality_condition}
%\end{equation}
%which states that the polarization of the electromagnetic wave is perpendicular to its propagation direction.

Any vector field $\vb A(\vb r,t)$ can be uniquely decomposed into a \emph{transverse} (divergence-free) component perpendicular to $\vb q$ and a
\emph{longitudinal} (curl-free) component parallel to $\vb q$.
In the Coulomb gauge framework, defined by the condition
\begin{equation}
  \nabla \cdot \vb A_{\mathrm{ext}}(\vb r,t) = 0,
\label{eq:coulomb_gauge_condition}
\end{equation}
the vector potential is purely transverse: it carries only the transverse part of the electric field and the magnetic field,
while the scalar potential $\Phi_{\mathrm{ext}}$ carries the longitudinal part of the electric field exclusively.
For a plane wave, this becomes apparent upon applying the
gauge condition~\eqref{eq:coulomb_gauge_condition} to
Eq.~\eqref{eq:plane_wave_vector_potential}, which yields
the transversality condition
\begin{equation}
  \vb q \cdot \vb A_0 = 0,
  \label{eq:transversality_condition}
\end{equation}
requiring that the polarization of the radiation field is perpendicular to its propagation direction.


%the longitudinal (parallel to the propagation direction vector $\vb q$) and  transverse (perpedicular to $\vb q$) components of  electric field
%are carried, respectively,   by the  a scalar and  vector potential exclusively.
%That is, the vector potential carries the transverse electric field and the magnetic field, 
% as can be seen by  a direct  application of Eq. \ref{eq:coulomb_gauge_condition}  to a Eq. \ref{eq:plane_wave_vector_potential}
% %resulting in the transversality condition
%\begin{equation}
%  \vb q \cdot \vb A_0 = 0,
%\label{eq:transversality_condition}
%\end{equation}

 In general,  the electromagnetic  potentials ($\vb A$ and $\Phi$)  and the KS potential all appear simultaneously. However, the gauge  freedom  allows  to transform  
 $\vb A$ and $\Phi$ to form suitable for a given problem using an arbitrarily selected scalar function $\chi(r,t)$  and  a set of  tranformation rules
\begin{align}
	& \vb A  \to \vb  A +\nabla\chi(r,t) , \\
	& \Phi \to \Phi-   \frac{1}{c}  \pdv{\chi(r,t}{t} ,
  \label{eq:gauge-rules}
%  \\
%  & \hat{U} = e^{-\imath\chi} = e^{\imath \vb r\cdot\vb A(t)} \label{eq:U-def} \\
%  & \Psi \to \Psi' = \hat{U}\Psi. 
\end{align}
%The Hamiltonian transforms as:
%\begin{equation}
%  \hat{H}
%  \to\hat{H}'
%  = \hat{U}\hat{H}\hat{U}^\dagger
%    - \imath \hat{U}\frac{\partial\hat{U}^\dagger}{\partial t}.
%  \label{eq:H-transform}
%\end{equation}
A specific choice of $\chi$ relevant to periodic systems is discussed in Sec.~\ref{sec:velocity_gauge}.

%In  chemistry application  the most common choice for RT-TDDFT  is so called  \emph{Length Gauge} 
% which  sets  $\vb A = 0$,  and   the external EM   field is introduced  trough a curl-free and magnetic-field free
%scalar potential
%\begin{equation}
%\Phi(r,t)= -\vb E \cdot \vb r,
%\label{eq:EM_dipolar}
%\end{equation}
%where  $\vb E=const$ represents a uniform (homogenuos) electric field,  and is usually used to investigate an optical response of finite size molecular system.
%The potential in   Eq. \ref{eq:EM_dipolar}  represents so  dipolar  scalar potential and can obtained  in long wavelength limit ($\vb q \to0$) of $e^{\imath \vq cdot\vbr}$
%in the  plane wave  Eq. \ref{eq:plane_wave_vector_potential}.
%This approximation is valid when size of molecular system $L$  is such that  $|\vb q| L<<1 $  meaning thatm olecualr size $L$ is  mcuh smaller than 
%the radiation wavelength  $\lambda =\frac{2 \pi}{|\vb q|}$, so the field is effectively uniform over the system. 

 
In chemistry applications, the most common choice for RT-TDDFT is the \emph{length gauge}, 
in which $\vb A = 0$ and the external electromagnetic field
enters entirely through a magnetic-field-free scalar potential.
In the long-wavelength limit, applicable when the size of the system $L$ satisfies $|\vb q| L \ll 1$ (i.e.\ $L$ is much
smaller than the radiation wavelength $\lambda = 2\pi/|\vb q|$), the external field is effectively
spatially uniform over the system and the interaction reduces to the electric-dipole approximation.
A spatially uniform electric field $\vb E(t)$ is represented by the dipolar scalar potential
\begin{equation}
  \Phi(\vb r,t) = -\vb E(t) \cdot \vb r,
  \label{eq:EM_dipolar}
\end{equation}
as can be verified from $\vb E = -\nabla\Phi$. This form is widely used to investigate the optical response of finite molecular systems.

While the plane-wave form~\eqref{eq:plane_wave_vector_potential} describes a monochromatic perturbation, in practice the
time-oscillating factor may be replaced by an arbitrary envelope function $f(t)$, so that the applied electric field
takes the form $\vb E(t) = \hat{\vb e}\, E_0\, f(t)$, where $\hat{\vb e}$ is the polarization direction.
A particularly useful choice is the Dirac delta kick, $f(t) = \kappa\,\delta(t)$, which delivers a broadband
impulsive perturbation of strength $\kappa$ at $t=0$.  Since the delta function has a flat spectral content,
a single time propagation yields the dynamical response at all frequencies simultaneously.
The frequency-dependent absorption spectrum is then obtained from the Fourier transform of the time-dependent observable
monitored during the propagation.

For finite systems, the relevant observable is the time-dependent dipole moment 
$  {\vb d}(t) = \int \vb r\rho(\vb r,t) d\vb r^3$, 
from which the polarizability and absorption cross section are obtained from  a dipole correlation function  by Fourier transform.
For periodic systems, however, the scalar-potential form $\Phi = -\vb E(t)\cdot\vb r$ is not compatible with
periodic boundary conditions, because the position operator $\vb r$ is discontinuous at the cell boundary.
This motivates the use of the velocity-gauge formulation, in which the external field enters through a spatially
uniform vector potential $\vb A(t)$ and the scalar potential is set to zero.
For finite systems, the optical response is extracted from the time-dependent dipole moment, which depends
on the density, whereas for truly periodic systems the relevant observable is the current density
$\vb j(\vb r,t)$, which remains well defined under periodic boundary conditions.



In a summary, it is useful at this point to recognize that external
perturbations in RT-TDDFT fall into two physically distinct categories.
The first category comprises electromagnetic field
perturbations, in which the external field enters through
the vector potential in the kinetic energy via the minimal
coupling $\vb{\hat p} \to \vb{\hat p} - (q_e/c)\vb A$.
The choice of gauge (length or velocity) determines how
this coupling is represented in the Hamiltonian, and the
long-wavelength limit leads to the dipolar approximation
discussed above.
Typical applications include optical absorption spectra
and laser-driven electron dynamics.

The second category comprises scalar potential
perturbations of the form
\begin{equation}
  \Phi_{\mathrm{ext}}(\vb r,t)
  = \Phi_0\, e^{\imath(\vb q\cdot\vb r - \omega t)}
    + \text{c.c.},
  \label{eq:scalar_probe}
\end{equation}
which are added directly to the Kohn--Sham potential
$\hat V_{\mathrm{KS}}$ and involve no gauge freedom.
The electric field generated by such a potential,
$\vb E = -\nabla\Phi_{\mathrm{ext}}
\propto \imath\,\vb q\, e^{\imath(\vb q\cdot\vb r - \omega t)}$,
is purely longitudinal (parallel to $\vb q$) and carries
no magnetic field.
This type of perturbation probes the density--density
response function $\chi(\vb q,\omega)$ and is the
natural tool for studying collective excitations at
finite momentum transfer, such as plasmon dispersion
and charge-density oscillations.

Although both categories may involve the same spatial
factor $e^{\imath\vb q\cdot\vb r}$, they are physically
and computationally distinct: the electromagnetic
perturbation couples to the electron momentum through
the kinetic energy and requires gauge treatment, while
the scalar perturbation couples to the electron density
through the potential energy and enters the calculation
without any gauge considerations.


Many important spectroscopies and collective excitations
involve finite momentum transfer and lie beyond the
dipolar approximation.
These include plasmon dispersion
$\varepsilon(\vb q,\omega)$, electron energy-loss
spectroscopy (EELS) where the beam transfers finite
$\vb q$ to the sample, inelastic X-ray scattering (IXS),
and magnon dispersion in magnetic systems.
Table~\ref{tab:gauge_cases} summarizes the common
external-field representations used in TDDFT and
their typical applications.

\begin{table}[t]
\footnotesize
\caption{Summary of common external-field representations
  used in TDDFT, applicable to both linear-response and
  real-time formulations.
  In Coulomb gauge, the scalar potential carries the
  longitudinal part of the electric field, while the
  transverse vector potential carries the transverse
 electric field and the magnetic field.}
\label{tab:gauge_cases}
\centering
\begin{tabular}{p{2.5cm}p{4.2cm}p{3.2cm}p{1.4cm}p{4.0cm}}
\hline
Case
  & Potentials
  & Field character
  & $B$ field
  & Typical applications \\
\hline
Length gauge, dipole
  & $\Phi=-\mathbf r\cdot\mathbf E(t)$,
    $\mathbf A=0$
  & homogeneous electric field
  & no
  & optical absorption of finite
    systems,
    photochemistry
    \cite{Yabana1996,Jakowski2025} \\[0.6em]
Velocity gauge
  & $\Phi=0$,
    $\mathbf A=\mathbf A(t)$
  & homogeneous field
  & no
  & optical response of solids and
    periodic systems under PBC
    \cite{Bertsch2000,Mattiat2022,Hanasaki2023}
    \\[0.6em]
Finite-$\mathbf q$ scalar probe
  & $\Phi \propto e^{i(\mathbf q\cdot\mathbf r
    -\omega t)}$, $\mathbf A=0$
  & longitudinal,
    $\mathbf E \parallel \mathbf q$
  & no
  & density response, loss function,
    EELS, plasmon
    dispersion
    \cite{Onida2002,Weissker2006,Weissker2010}
    \\[0.6em]
Finite-$\mathbf q$ EM wave
  & $\mathbf A \propto
    \boldsymbol{\epsilon}\,
    e^{i(\mathbf q\cdot\mathbf r-\omega t)}$,
    $\mathbf q\cdot\boldsymbol{\epsilon}=0$
  & transverse EM radiation
  & yes
  & IXS, radiative coupling
    beyond dipole,
    magnon dispersion
    \cite{TancogneDejean2020,Krieger2015}
    \\
\hline
\end{tabular}
\end{table}

%%%%%%%%%%%%%%%%%%%%%%%%%%%%%%%%%%%%%%%%%%%%%%%%%%%%%%%%%%%%%%%%%%%%%%%%%%%%%%%%%%
\subsection{ VELOCITY GAUGE}


For finite systems, the relevant observable is the
time-dependent dipole moment $\vb d(t) = \int \vb r\,
\rho(\vb r,t) d\vb r$, from which the polarizability
and absorption cross section follow by Fourier transform.
For periodic systems, however, the scalar-potential form
$\Phi = -\vb E(t)\cdot\vb r$ is not compatible with
periodic boundary conditions, because the position
operator $\vb r$ is discontinuous at the cell boundary.
This motivates the use of the velocity-gauge formulation,
in which the external field enters through a spatially
uniform vector potential $\vb A(t)$ and the scalar
potential is set to zero.

The time-dependent Kohn--Sham equation then takes the form
\begin{equation}
  \imath \pdv{}{t} \Psi_j(\vb r,t) = \left[
    \frac{1}{2m}\left(\vb{\hat p}
    - \frac{q_e}{c}\vb A(t)\right)^2
    + \hat V_{\mathrm{KS}}[\rho](t)
  \right] \Psi_j(\vb r,t),
  \label{eq:tdks_velocity_gauge}
\end{equation}
where the two gauges are related by the Göppert-Mayer
transformation, discussed in detail in
Sec.~\ref{sec:velocity_gauge}.
In the velocity gauge, the position operator no longer
appears in the Hamiltonian and all terms are compatible
with periodic boundary conditions.
Correspondingly, the physical observable used to extract
the optical response is the time-dependent current density
$\vb j(\vb r,t)$ rather than the dipole moment, since
the current is well defined under periodic boundary
conditions whereas $\vb d(t)$ is not.





